% ============================================================
% Trabajo Especial de Licenciatura en Ciencias de la Computación
% FAMAF - Universidad Nacional de Córdoba
% ============================================================

\documentclass[12pt,a4paper]{report}


\usepackage{fontspec}
\usepackage{polyglossia}
\setdefaultlanguage{spanish}

\IfFontExistsTF{Times New Roman}
  {\setmainfont{Times New Roman}}
  {\setmainfont{TeX Gyre Termes}}

\usepackage[
    a4paper,
    top=2.5cm,
    bottom=2.5cm,
    left=3cm,
    right=2.5cm
]{geometry}

\usepackage{setspace}
\setstretch{1.0}

\setlength{\parindent}{1.25cm}
\setlength{\parskip}{0pt}

\usepackage{amsmath,amssymb}
\usepackage{graphicx}
\usepackage{booktabs}
\usepackage{float}
\usepackage{caption}
\usepackage{subcaption}

\usepackage[
    hidelinks,
    bookmarks
]{hyperref}

\usepackage[
    capitalise,
    nameinlink
]{cleveref}

\usepackage[
    backend=biber,
    style=alphabetic,
    language=spanish,
    autolang=hyphen
]{biblatex}

\addbibresource{bibliografia.bib}

\usepackage{fancyhdr}

\pagestyle{fancy}
\fancyhf{}
\fancyhead[R]{\thepage}
\renewcommand{\headrulewidth}{0pt}

\fancypagestyle{plain}{
    \fancyhf{}
    \fancyfoot[C]{\thepage}
    \renewcommand{\headrulewidth}{0pt}
}

\usepackage{titlesec}

\titleformat{\chapter}[display]
    {\normalfont\huge\bfseries}
    {Capítulo \thechapter}
    {10pt}
    {\Huge}

\titlespacing*{\chapter}
    {0pt}
    {0pt}
    {25pt}

% ------------------------------------------------------------
% Datos del trabajo
% ------------------------------------------------------------
\newcommand{\tituloTrabajo}{Computación óptica}
\newcommand{\autorTrabajo}{Tomás Achával Berzero}
\newcommand{\directorTrabajo}{Dr. Federico Pont}
\newcommand{\codirectorTrabajo}{Dr. Ing. James Kunst}
\newcommand{\fechaTrabajo}{\today}


\begin{document}

% HOJA EN BLANCO INICIAL — guarda
\thispagestyle{empty}
\mbox{}
\clearpage

% CARÁTULA
\thispagestyle{empty}

\begin{center}

    % Logos
    \begin{minipage}[t]{\textwidth}
        \raggedright
        \IfFileExists{logos/unc.png}
            {\includegraphics[height=1.7cm]{logos/unc.png}}
            {\fbox{\rule{0pt}{1.5cm}\rule{3cm}{0pt}}}
        \hfill
        \IfFileExists{logos/famaf.png}
            {\includegraphics[height=1.7cm]{logos/famaf.png}}
            {\fbox{\rule{0pt}{1.5cm}\rule{3cm}{0pt}}}
    \end{minipage}

    \vspace{3cm}

    {\Huge\bfseries \tituloTrabajo\par}

    \vspace{0.7cm}

    {\Large Trabajo Especial de Licenciatura\par}

    \vspace{1.5cm}

    por

    \vspace{0.3cm}

    {\Large \autorTrabajo\par}

    \vspace{2cm}

    Presentado ante la Facultad de Matemática, Astronomía, Física y
    Computación como parte de los requerimientos para la obtención del
    grado de Licenciado en Ciencias de la Computación de la

    \vspace{0.4cm}

    {\large\bfseries UNIVERSIDAD NACIONAL DE CÓRDOBA\par}

    \vspace{1.5cm}

    {\large \fechaTrabajo\par}

    \vfill

    \begin{large}
        \textbf{Director:} \directorTrabajo

        \vspace{0.3cm}

        \textbf{Co-Director:} \codirectorTrabajo
    \end{large}

\end{center}

\clearpage

% RESUMEN
\chapter*{Resumen}
\addcontentsline{toc}{chapter}{Resumen}

% Escribir aquí el resumen en castellano.

\vspace{1cm}

\noindent\textbf{Palabras clave:}

\clearpage

% ABSTRACT
\chapter*{Abstract}
\addcontentsline{toc}{chapter}{Abstract}

% Write the abstract in English here.

\vspace{1cm}

\noindent\textbf{Keywords:}

\clearpage


% AGRADECIMIENTOS
\chapter*{Agradecimientos}
\addcontentsline{toc}{chapter}{Agradecimientos}

% Opcional.

\clearpage

% ÍNDICE
\tableofcontents
\clearpage

% CONTENIDO
\chapter{Introducción}

\section{Motivación}

\section{Objetivos}

\section{Organización del trabajo}


\chapter{Orígenes de la computación óptica}

Invención del láser 1960

\section{Primeras propuestas}

Señales de radar

DOC-II

Pattern recognition

Memoria Optica (CDs)

\section{Computación óptica digital}

Más detalles de la DOC-II

\section{Consolidación de la electrónica}

Ley de Moore

Avances en los procesos de fabricación

Problemas con el tamaño y accesibilidad de los dispositivos


\chapter{Fundamentos físicos}

En este capítulo se presentarán los fundamentos físicos que permiten maniplar luz con luz, habilitando el paradigma de computación all-optical.

\section{Formas de describir la luz}

\section{Propagación e Interferencia}

\subsection{No linealidades}

\section{Dispositivos ópticos básicos}

Delay Lines.
Mach-Zehnder.
Memoria óptica (CDs).


\chapter{Modelado y simulación}
Para comprender en profundidad el funcionamiento de los dispositivos ópticos, se llevó a cabo la simulación de un transistor all-optical basado en un interferómetro Mach-Zehnder que contiene una no linealidad óptica de tipo III en una de sus ramas.

\section{Proceso de modelado}

\section{Metodología de simulación}

\section{Circuitos ópticos}

\section{Resultados}



\chapter{Resurgimiento de la computación óptica}

\section{Limitaciones de la electrónica}

Frecuencia del clock.

Temperatura.

Ley de moore actual.

\section{Materiales y procesos de fabricación}

\subsection{Primeros materiales}
\subsection{Silicon photonics}

\section{El avance de las interconexiones ópticas}

\section{Eficiencia energética}



\chapter{Arquitecturas ópticas modernas}

\section{Aceleradores ópticos}

\section{Computación all-optical}

\subsection{Akhetonics}

Subsección destinada a la entrevista con kissner (idealmente a la segunda, en enero/febrero).

\subsection{Principales desafíos técnicos}

\subsection{Realidad del consumo energético}

\subsection{Escalabilidad}

\section{Perspectivas futuras}


\chapter{Conclusiones}
Que esperar en los próximos años?


% BIBLIOGRAFÍA
\printbibliography[heading=bibintoc,title={Bibliografía}]

% ANEXOS
\appendix

\chapter{Material complementario}

% HOJA FINAL DE APROBACIÓN
\clearpage
\thispagestyle{empty}

\vspace*{\fill}

\begin{center}
    \Large\textbf{Constancia de aprobación}
\end{center}

\vspace{1cm}

% Hoja destinada a la constancia de correspondencia entre
% la versión presentada y la versión aprobada por el tribunal.

\vspace{3cm}

\noindent
\begin{minipage}{0.45\textwidth}
    \centering
    \rule{6cm}{0.4pt}\\
    Firma
\end{minipage}
\hfill
\begin{minipage}{0.45\textwidth}
    \centering
    \rule{6cm}{0.4pt}\\
    Firma
\end{minipage}

\vspace*{\fill}

% HOJA EN BLANCO FINAL — guarda
\clearpage
\thispagestyle{empty}
\mbox{}

\end{document}
